\scanpage{54}
de correspondance divisorielle canonique
\[
c'\in\mathrm{Cl}(B\unsure{,|\times}\deleted{\unsure{A_Y|A}}Y).
\]
\deleted{\unsure{Nous prop.|Nous supposons}} Cette classe définit un homomorphisme \deleted{$Y\longrightarrow\underline{\mathrm{Pic}}^{\unsure{1|0}}_{B/k}$, d'où un homomorphisme} $A_Y\longrightarrow\underline{\mathrm{Pic}}^0_{B/k}$, or par Cartier (bidualité) on a \inserted{un isom can.} $\underline{\mathrm{Pic}}^0_{B/k}\simeq A_X$, donc $c'$ définit un \deleted{\unsure{élément}} homomorphisme de VA
\[
u_Z:A_Y\longrightarrow A_X,
\]
ne dépendant que de la classe d'équivalence numérique $z$ de $Z$. Donc le th. 1 résultera de l'énoncé plus précis :

\uline{Corollaire} 1.3. L'homomorphisme $u_Z:A_Y\unsure{-|\longrightarrow} A_X$ \unsure{qui on|qu'on} vient de construire satisfait à la condition envisagée dans le th. 1.1. \deleted{\unsure{; en|car}}

Cela résultera du

\uline{Lemme} 1.4. Soit $\unsure{\varpi|\mathfrak{a}|\mathfrak{A}}$ un diviseur de degré $0$
\par\noindent\begin{tabular}{@{}|p{.92\linewidth}@{}}
sur $Y$ tel que $ {}^{\deleted{\unsure{t|*}}}Z(\varpi)$ soit défini. Alors l'élément $\Theta(S_X(Z(\varpi)))$ de $\mathrm{Pic}^0(A_X)$ est égal à $c'(\varpi)$.
\end{tabular}
% archive sheet number: 53
